\section{Các nghiên cứu liên quan}
Sự phát triển của NIDS hiện đại xoay quanh ba bài toán cốt lõi: xử lý mất cân bằng dữ liệu, nâng cao độ chính xác phát hiện, và tối ưu triển khai tại vùng biên với tài nguyên hạn chế. Phần này điểm lại các hướng tiếp cận tiêu biểu theo ba trục đó.

\subsection{Kiến trúc học sâu lai và tối ưu hóa vùng biên}
Nhằm nâng cao độ chính xác, nhiều nghiên cứu kết hợp các kiến trúc học sâu như mạng tích chập (CNN), bộ nhớ dài-ngắn hạn (LSTM/GRU) và cơ chế chú ý (Attention)~\cite{hu2024improved, ji2025hybrid} để trích xuất đồng thời đặc trưng không gian và thời gian của lưu lượng mạng. Mặc dù đạt độ chính xác cao, các mô hình này thường có số tham số lớn, độ trễ suy luận cao và tiêu thụ nhiều năng lượng, gây khó khăn khi triển khai trên thiết bị biên. Để khắc phục, xu hướng nghiên cứu đang dịch chuyển sang thiết kế mạng tinh gọn: các giải pháp như Autoencoder--LSTM~\cite{hejazi2025lightweight_botnet_iot} hay mạng học sâu lai tối giản~\cite{altaie2024hybrid_lightweight_ids} đã giảm đáng kể chi phí vận hành và được kiểm chứng trên nền tảng nhúng như Raspberry Pi. Tuy nhiên, việc nén mô hình theo cách này thường đánh đổi một phần độ chính xác, đặc biệt với các lớp tấn công hiếm. Ở bối cảnh trong nước, hướng phát hiện tấn công IoT bằng học máy cũng được quan tâm: Thái và cộng sự~\cite{thai2022twotier} đề xuất kiến trúc IDS hai tầng (mô hình nhẹ tại gateway, mô hình phức tạp hơn trên máy chủ), còn Vũ và Nguyễn~\cite{quang2025featurereduction} khảo sát các kỹ thuật giảm chiều vector đặc trưng nhằm tăng hiệu quả tính toán cho mô hình phát hiện. Tuy vậy, các nghiên cứu này chưa khai thác chưng cất tri thức để nén mô hình tổng hợp xuống mức triển khai trực tiếp tại biên.

\subsection{Tái cân bằng dữ liệu và phân tích thời gian thực}
Để duy trì tốc độ phân tích trên luồng dữ liệu lớn, các hệ thống NIDS thường tận dụng nền tảng xử lý dòng phân tán như Apache Kafka và Spark Streaming~\cite{kumar2025kafka_spark_ids}. Song song, vấn đề mất cân bằng lớp---vốn làm suy giảm nghiêm trọng khả năng phát hiện tấn công hiếm---được giải quyết theo hai nhánh chính. Nhánh sinh dữ liệu sử dụng các mô hình đối kháng tạo sinh tiên tiến như TMG-GAN~\cite{ding2024tmg, shamim2025gan_sdn_ids} để tái tạo mẫu tấn công thiểu số, song chi phí huấn luyện cao và tiềm ẩn nguy cơ sinh ra mẫu thiếu thực tế. Nhánh lấy mẫu (sampling) gồm các kỹ thuật tăng/giảm mẫu; trong đó việc kết hợp giảm mẫu lớp đa số với lọc nhiễu ranh giới (như ENN) giúp tạo ranh giới phân lớp sắc nét với chi phí tính toán thấp hơn---là hướng mà nghiên cứu này kế thừa.

\subsection{Học tổng hợp và chưng cất tri thức cho NIDS}
Học tổng hợp (ensemble), đặc biệt là cơ chế Stacking~\cite{wolpert1992stacked}, đã chứng tỏ hiệu quả vượt trội trong phát hiện xâm nhập nhờ kết hợp điểm mạnh của nhiều bộ phân loại cơ sở. Trên CICIDS2017, các kiến trúc stacking---từ tổ hợp học sâu CNN+TCN+LSTM~\cite{amara2025stacked} đến stacking kết hợp lựa chọn đặc trưng lai~\cite{huang2025stacking, jouybari2025stacking}---đều đạt độ chính xác rất cao và tỷ lệ cảnh báo giả thấp. Tuy vậy, các mô hình tổng hợp này có kích thước và độ trễ lớn, khó triển khai trực tiếp tại biên. Chưng cất tri thức (Knowledge Distillation)~\cite{hinton2015distilling} là một hướng nén mô hình hứa hẹn, chuyển giao tri thức từ mô hình ``Giáo viên'' phức tạp sang mô hình ``Học sinh'' gọn nhẹ; tuy nhiên phần lớn nghiên cứu hiện tập trung cho mạng nơ-ron, ít khai thác cho các bộ phân loại dạng cây---vốn đặc biệt phù hợp với thiết bị biên nhờ tốc độ suy luận nhanh và kích thước nhỏ.

Mặc dù đã có nhiều bước tiến, các nghiên cứu hiện tại vẫn chưa giải quyết trọn vẹn bài toán tối ưu \emph{đồng thời} ba yếu tố: nâng cao độ chính xác, xử lý hiệu quả mất cân bằng dữ liệu, và đáp ứng yêu cầu độ trễ thấp của điện toán biên. Kiến trúc NIDS hai giai đoạn đề xuất trong nghiên cứu này---tích hợp chiến lược lấy mẫu lai RUS-ENN, mô hình tổng hợp Stacking và chưng cất tri thức nhãn mềm sang cây quyết định---được thiết kế nhằm khắc phục đồng thời các hạn chế nêu trên.
